\documentclass[11pt,a4paper]{article}

% -----------------------------------------------------------------------------
% CS6868 Research Mini-Project Report Template
% -----------------------------------------------------------------------------
% Target length: 5-10 pages (including figures, excluding references).
% Compile with: latexmk -pdf main.tex
% -----------------------------------------------------------------------------

\usepackage[utf8]{inputenc}
\usepackage[T1]{fontenc}
\usepackage[margin=1in]{geometry}
\usepackage{microtype}
\usepackage{lmodern}
\usepackage{amsmath,amssymb}
\usepackage{graphicx}
\usepackage{booktabs}
\usepackage{array}
\usepackage{enumitem}
\usepackage{xcolor}
\usepackage{listings}
\usepackage{hyperref}
\usepackage[capitalise,noabbrev]{cleveref}

\hypersetup{
  colorlinks=true,
  linkcolor=black,
  citecolor=blue!60!black,
  urlcolor=blue!60!black,
}

% --- Listings style (for code snippets) ---------------------------------------
\definecolor{codegray}{gray}{0.95}
\lstset{
  basicstyle=\ttfamily\small,
  backgroundcolor=\color{codegray},
  frame=single,
  framesep=4pt,
  breaklines=true,
  columns=fullflexible,
  keepspaces=true,
  showstringspaces=false,
}

% --- Title block --------------------------------------------------------------
\title{%
  \textbf{CS6868 Research Mini-Project} \\[0.3em]
  \Large Your Project Title Here%
}
\author{%
  Member One \\ \texttt{rollno1@smail.iitm.ac.in}
  \and
  Member Two \\ \texttt{rollno2@smail.iitm.ac.in}
  \and
  Member Three \\ \texttt{rollno3@smail.iitm.ac.in}
}
\date{\today}

\begin{document}
\maketitle

\begin{abstract}
Briefly (150--200 words) describe the problem you tackled, your approach, the
headline evaluation results, and the main conclusion. Readers (the graders)
should understand the essence of your project from the abstract alone.
\end{abstract}

% =============================================================================
\section{Goals}
\label{sec:goals}
% =============================================================================

State the goals of the project. What concurrent algorithm or data structure
did you implement? Why is it interesting? How does it connect to course
topics (cite the relevant lecture(s) from the course)?

End this section with a crisp statement of your \emph{research question}
(borrowed from or adapted from the chosen project idea in
\texttt{project\_ideas.md}, or formulated by you if you proposed your own
topic).

% =============================================================================
\section{Background}
\label{sec:background}
% =============================================================================

Summarise the algorithm or technique at a level that a classmate who has
attended the lectures but not read the paper can follow. Cite the primary
references~\cite{example}. If helpful, include a small figure or
pseudo-code listing.

% Example listing:
% \begin{lstlisting}[language=caml,caption={Sketch of the algorithm.},label={lst:algo}]
% let foo x = ...
% \end{lstlisting}

% =============================================================================
\section{Tasks Undertaken}
\label{sec:tasks}
% =============================================================================

Describe what you actually built. Organise this section as sub-sections per
major task; use the \emph{Tasks} list from the project idea as a checklist.

\subsection{Implementation}
Describe the implementation: key data structures, atomic operations used,
subtle invariants, and any design choices you had to make. Link to the
relevant files/commits in your GitHub repository.

\subsection{Testing and verification}
Describe how you tested correctness (QCheck-Lin, QCheck-STM, dscheck, TSAN,
etc.) and what bugs the tests caught.

% =============================================================================
\section{Evaluation}
\label{sec:evaluation}
% =============================================================================

Describe your experimental setup (hardware, number of domains/threads, how
you measured, how many runs, how you handled warm-up and noise). Present
results with plots or tables.

% Example figure:
% \begin{figure}[h]
%   \centering
%   \includegraphics[width=0.8\linewidth]{figures/throughput.pdf}
%   \caption{Throughput vs. threads under the balanced workload.}
%   \label{fig:throughput}
% \end{figure}

Discuss what the numbers \emph{mean}. Where did the algorithm scale well?
Where did it break down? Did the results match your prior expectations?

% =============================================================================
\section{Reflection on the Use of LLMs}
\label{sec:llm}
% =============================================================================

% This section is required and is graded as part of the report. Be specific
% and honest -- we are more interested in what you learned from working with
% the LLM than in a polished narrative.

\paragraph{Tools and models used.}
List every LLM-backed tool you used (e.g., GitHub Copilot, ChatGPT, Claude,
Gemini, Cursor), the specific model version where relevant
(e.g., \emph{GPT-5}, \emph{Claude Sonnet 4.6}, \emph{Gemini 2.5 Pro}), and
roughly for what --- code generation, debugging, explaining a paper, writing
test cases, drafting prose, generating plots, etc.

\paragraph{What worked well.}
Concrete examples where the LLM genuinely accelerated your work. Did it
correctly implement a subtle CAS loop on the first try? Did it explain a
paper better than the paper did? Did it catch a race you had missed? Cite
specific commits or passages if you can.

\paragraph{What did not work.}
Concrete examples where the LLM got it wrong or led you astray --- wrong
memory ordering, subtly broken retry logic, confident but incorrect
explanations, hallucinated API calls, plausible-looking citations that did
not exist, etc. How did you catch the mistake?

\paragraph{What was surprising.}
Anything you did not expect --- both positive (unexpectedly good at X) and
negative (unexpectedly bad at Y). Differences between models or tools?

\paragraph{What was difficult.}
Parts of the project where the LLM was \emph{unable to substitute for
understanding} --- reasoning about memory models, debugging a rare
interleaving, justifying linearisation points, choosing between design
alternatives. What did you end up having to learn yourself?

\paragraph{Your overall assessment.}
One or two sentences: did the LLM net help or hurt this project? Would you
use it the same way again?

% =============================================================================
\section{Conclusions}
\label{sec:conclusions}
% =============================================================================

Answer the research question stated in \cref{sec:goals}. Summarise the main
findings. State the limitations of your work honestly. Suggest what you
would do next given more time.

% =============================================================================
\section*{Contributions}
\label{sec:contributions}
% =============================================================================

% Required. Percentages must sum to 100. Every member should contribute.

\begin{center}
\begin{tabular}{@{}lcl@{}}
\toprule
\textbf{Member} & \textbf{Contribution (\%)} & \textbf{Main responsibilities} \\
\midrule
Member One   & 34\% & e.g., lock-free implementation, QCheck-Lin tests \\
Member Two   & 33\% & e.g., baseline implementations, benchmarking harness \\
Member Three & 33\% & e.g., evaluation plots, report, literature survey \\
\midrule
\textbf{Total} & \textbf{100\%} & \\
\bottomrule
\end{tabular}
\end{center}

% =============================================================================
\bibliographystyle{plainurl}
\bibliography{references}
% =============================================================================

\end{document}
